\documentclass[10pt,a4paper]{amsart}
\usepackage{amsmath,amssymb,amsthm,graphicx}
\usepackage{verbatim}
\newenvironment{reference}{\comment}{\endcomment}
%\newenvironment{reference}{}{}
\newenvironment{slogan}{\comment}{\endcomment}
\newenvironment{history}{\comment}{\endcomment}
% For commutative diagrams we use Xy-pic
\usepackage[all]{xy}

% We use 2cell for 2-commutative diagrams.
\xyoption{2cell}
% Theorem environments.
%
\theoremstyle{plain}
\newtheorem{theorem}[subsection]{Theorem}
\newtheorem{proposition}[subsection]{Proposition}
\newtheorem{lemma}[subsection]{Lemma}

\theoremstyle{definition}
\newtheorem{definition}[subsection]{Definition}
\newtheorem{example}[subsection]{Example}
\newtheorem{exercise}[subsection]{Exercise}
\newtheorem{situation}[subsection]{Situation}

\theoremstyle{remark}
\newtheorem{remark}[subsection]{Remark}
\newtheorem{remarks}[subsection]{Remarks}

\numberwithin{equation}{subsection}

% Macros
%
\def\lim{\mathop{\mathrm{lim}}\nolimits}
\def\colim{\mathop{\mathrm{colim}}\nolimits}
\def\Spec{\mathop{\mathrm{Spec}}}
\def\Hom{\mathop{\mathrm{Hom}}\nolimits}
\def\Ext{\mathop{\mathrm{Ext}}\nolimits}
\def\SheafHom{\mathop{\mathcal{H}\!\mathit{om}}\nolimits}
\def\SheafExt{\mathop{\mathcal{E}\!\mathit{xt}}\nolimits}
\def\Sch{\mathit{Sch}}
\def\Mor{\mathop{\mathrm{Mor}}\nolimits}
\def\Ob{\mathop{\mathrm{Ob}}\nolimits}
\def\Sh{\mathop{\mathit{Sh}}\nolimits}
\def\NL{\mathop{N\!L}\nolimits}
\def\CH{\mathop{\mathrm{CH}}\nolimits}
\def\proetale{{pro\text{-}\acute{e}tale}}
\def\etale{{\acute{e}tale}}
\def\QCoh{\mathit{QCoh}}
\def\Ker{\mathop{\mathrm{Ker}}}
\def\Im{\mathop{\mathrm{Im}}}
\def\Coker{\mathop{\mathrm{Coker}}}
\def\Coim{\mathop{\mathrm{Coim}}}

% Boxtimes
%
\DeclareMathSymbol{\boxtimes}{\mathbin}{AMSa}{"02}

%
% Macros for moduli stacks/spaces
%
\def\QCohstack{\mathcal{QC}\!\mathit{oh}}
\def\Cohstack{\mathcal{C}\!\mathit{oh}}
\def\Spacesstack{\mathcal{S}\!\mathit{paces}}
\def\Quotfunctor{\mathrm{Quot}}
\def\Hilbfunctor{\mathrm{Hilb}}
\def\Curvesstack{\mathcal{C}\!\mathit{urves}}
\def\Polarizedstack{\mathcal{P}\!\mathit{olarized}}
\def\Complexesstack{\mathcal{C}\!\mathit{omplexes}}
% \Pic is the operator that assigns to X its picard group, usage \Pic(X)
% \Picardstack_{X/B} denotes the Picard stack of X over B
% \Picardfunctor_{X/B} denotes the Picard functor of X over B
\def\Pic{\mathop{\mathrm{Pic}}\nolimits}
\def\Picardstack{\mathcal{P}\!\mathit{ic}}
\def\Picardfunctor{\mathrm{Pic}}
\def\Deformationcategory{\mathcal{D}\!\mathit{ef}}
\usepackage{cleveref}
\crefname{originallemma}{original lemma}{original lemmas}
\Crefname{originallemma}{Original lemma}{Original lemmas}
\makeatletter
\newlabel{lemma-properties-diagonal}{{lemma-properties-diagonal}{source}{}{sourcepointer-1}{}}
\newlabel{lemma-properties-diagonal@cref}{{[originallemma][0][]lemma-properties-diagonal}{[1][0][]source}{}{}{}}
\makeatother
\crefname{originalreference}{original reference}{original references}
\Crefname{originalreference}{Original reference}{Original references}
\makeatletter
\newlabel{stacks-properties-lemma-points-cartesian}{{stacks-properties-lemma-points-cartesian}{source}{}{sourcepointer-2}{}}
\newlabel{stacks-properties-lemma-points-cartesian@cref}{{[originalreference][0][]stacks-properties-lemma-points-cartesian}{[1][0][]source}{}{}{}}
\makeatother
\crefname{originalreference}{original reference}{original references}
\Crefname{originalreference}{Original reference}{Original references}
\makeatletter
\newlabel{stacks-properties-section-surjective}{{stacks-properties-section-surjective}{source}{}{sourcepointer-3}{}}
\newlabel{stacks-properties-section-surjective@cref}{{[originalreference][0][]stacks-properties-section-surjective}{[1][0][]source}{}{}{}}
\makeatother
\title{A morphism is universally injective exactly when its diagonal is surjective}
\author{The Stacks Project Authors}
\date{stacks-project:stacks-morphisms.tex:local-snapshot-20261012; GNU FDL 1.2 or later}
\begin{document}
\maketitle
\noindent\textit{Editorial note.} All delivered mathematical text is the author's own, copied from the pinned Stacks Project chapter source sources/stacks-algebraic-stacks-morphisms/pinned-original-source.tex; no statement or proof step was rewritten, reordered or removed.\par\medskip
\noindent\textit{Editorial note.} Preamble provenance: the Stacks Project keeps one single shared preamble for every chapter; this card takes it directly from the staged sources/stacks-algebraic-stacks-morphisms/preamble.tex as three verbatim spans (lines 8-12, 14-18 and 154-217), so no derived file was written and nothing was added or redefined. The wrapper class and the base packages are supplied by the assembler.\par\medskip
\noindent\textit{Editorial note.} Bibliography provenance: neither the counted statement nor its proof contains a citation (this lane rescans the whole span before writing the card), so the card carries no bibliography and no entry was reconstructed.\par\medskip
\noindent\textit{Editorial note.} Selection provenance: this card was opened under the file-based control contract (CONTROL.md), with the mechanical pre-card gate run on its statement and proof interval before the card was written.\par\medskip
\noindent\textit{Allowed standard prerequisites (editorial).}\begin{itemize}
\item The factorization of the diagonal through the graph of the morphism and the base change of the diagonal by the source. Used inside the counted proof (source lines 2923-2933). Identifies surjectivity of the diagonal with surjectivity of its pullback along the identity section.
\item The criterion that a morphism of algebraic spaces is surjective if and only if it is surjective on points after base change to fields. Invoked inside the counted proof (source lines 2934-2939) as a cross-chapter pointer. Reduces surjectivity of the diagonal to a condition on points over fields.
\item The description of the points of a fibre product of stacks over a field as pairs of points isomorphic after base change. Used inside the counted proof (source lines 2940-2960). Turns the pointwise condition into the statement that two points of the source with the same image are equal, which is universal injectivity.
\end{itemize}\medskip
\ifdefined\hypertarget\hypertarget{sourcepointer-1}{}\fi
\noindent\textit{Source reference.} Original lemma lemma-properties-diagonal: sources/stacks-algebraic-stacks-morphisms/pinned-original-source.tex, source lines 420--420. Navigation pointer only: the counted text refers to the Stacks item lemma-properties-diagonal, which belongs to another chapter of the project and is not part of this card. The pointer carries the target's own label as its number because this chapter's staged snapshot does not contain the other chapter's tag table; its content is not reproduced here.\par\medskip
\ifdefined\hypertarget\hypertarget{sourcepointer-2}{}\fi
\noindent\textit{Source reference.} Original reference stacks-properties-lemma-points-cartesian: sources/stacks-algebraic-stacks-morphisms/pinned-original-source.tex, source lines 2126--2126. Navigation pointer only: the counted text refers to the Stacks item stacks-properties-lemma-points-cartesian, which belongs to another chapter of the project and is not part of this card. The pointer carries the target's own label as its number because this chapter's staged snapshot does not contain the other chapter's tag table; its content is not reproduced here.\par\medskip
\ifdefined\hypertarget\hypertarget{sourcepointer-3}{}\fi
\noindent\textit{Source reference.} Original reference stacks-properties-section-surjective: sources/stacks-algebraic-stacks-morphisms/pinned-original-source.tex, source lines 2931--2931. Navigation pointer only: the counted text refers to the Stacks item stacks-properties-section-surjective, which belongs to another chapter of the project and is not part of this card. The pointer carries the target's own label as its number because this chapter's staged snapshot does not contain the other chapter's tag table; its content is not reproduced here.\par\medskip
\begin{lemma}
\label{lemma-universally-injective}
Let $f : \mathcal{X} \to \mathcal{Y}$ be a morphism of algebraic stacks.
The following are equivalent
\begin{enumerate}
\item $f$ is universally injective,
\item $\Delta : \mathcal{X} \to \mathcal{X} \times_\mathcal{Y} \mathcal{X}$
is surjective, and
\item for an algebraically closed field, for
$x_1, x_2 : \Spec(k) \to \mathcal{X}$, and for a $2$-arrow
$\beta : f \circ x_1 \to f \circ x_2$ there is a
$2$-arrow $\alpha : x_1 \to x_2$ with
$\beta = \text{id}_f \star \alpha$.
\end{enumerate}
\end{lemma}

\begin{proof}
(1) $\Rightarrow$ (2). If $f$ is universally injective, then the
first projection
$|\mathcal{X} \times_\mathcal{Y} \mathcal{X}| \to |\mathcal{X}|$
is injective, which implies that $|\Delta|$ is surjective.

\medskip\noindent
(2) $\Rightarrow$ (1). Assume $\Delta$ is surjective. Then any base change
of $\Delta$ is surjective (see Properties of Stacks, Section
\ref{stacks-properties-section-surjective}).
Since the diagonal of a base change
of $f$ is a base change of $\Delta$, we see that it suffices
to show that $|\mathcal{X}| \to |\mathcal{Y}|$ is injective.
If not, then by Properties of Stacks, Lemma
\ref{stacks-properties-lemma-points-cartesian}
we find that the first projection
$|\mathcal{X} \times_\mathcal{Y} \mathcal{X}| \to |\mathcal{X}|$
is not injective. Of course this means that $|\Delta|$ is not
surjective.

\medskip\noindent
(3) $\Rightarrow$ (2). Let
$t \in |\mathcal{X} \times_\mathcal{Y} \mathcal{X}|$.
Then we can represent $t$ by a morphism
$t : \Spec(k) \to \mathcal{X} \times_\mathcal{Y} \mathcal{X}$
with $k$ an algebraically closed field.
By our construction of $2$-fibre products we can represent
$t$ by $(x_1, x_2, \beta)$ where $x_1, x_2 : \Spec(k) \to \mathcal{X}$
and $\beta : f \circ x_1 \to f \circ x_2$ is a $2$-morphism.
Then (3) implies that there is a $2$-morphism
$\alpha : x_1 \to x_2$ mapping to $\beta$.
This exactly means that $\Delta(x_1) = (x_1, x_1, \text{id})$
is isomorphic to $t$. Hence (2) holds.

\medskip\noindent
(2) $\Rightarrow$ (3). Let $x_1, x_2 : \Spec(k) \to \mathcal{X}$
be morphisms with $k$ an algebraically closed field. Let
$\beta : f \circ x_1 \to f \circ x_2$ be a $2$-morphism.
As in the previous paragraph, we obtain a morphism
$t = (x_1, x_2, \beta) :
\Spec(k) \to \mathcal{X} \times_\mathcal{Y} \mathcal{X}$.
By Lemma \ref{lemma-properties-diagonal}
$$
T = \mathcal{X}
\times_{\Delta, \mathcal{X} \times_\mathcal{Y} \mathcal{X}, t} \Spec(k)
$$
is an algebraic space locally of finite type over $\Spec(k)$.
Condition (2) implies that $T$ is nonempty. Then since $k$ is
algebraically closed, there is a $k$-point in $T$.
Unwinding the definitions this means there is a morphism
$\alpha : x_1 \to x_2$ in $\Mor(\Spec(k), \mathcal{X})$
such that $\beta = \text{id}_f \star \alpha$.
\end{proof}

\end{document}
